\documentclass[10pt,a4paper]{amsart}
\usepackage[margin=19mm]{geometry}
\usepackage{amsmath,amssymb,mathtools}
\usepackage[scr=rsfs,cal=euler]{mathalfa}
\usepackage{xfrac}
\usepackage[unicode,hidelinks,bookmarks=false]{hyperref}
\newcommand{\ie}{i.e.\ }
\newcommand{\cf}{cf.\ }
\newcommand{\resp}{respectively\ }
\newcommand{\Subsection}{\subsection}
\providecommand{\nobreakdash}{\nobreak\hbox{-}}
\setlength{\emergencystretch}{2em}
\multlinegap0pt
\usepackage{bm}
\usepackage{bbm}
\usepackage{dsfont}
\usepackage{xspace}
\usepackage{cancel}
\usepackage{mathtools}
\usepackage{amsthm}
\makeatletter
\@ifundefined{lemma}{\newtheorem{lemma}{Lemma}}{}
\makeatother
\let\tilde\widetilde

\setlength{\parskip}{0.5\baselineskip}%
\title[Operator limit of Wigner matrices I]{Operator limit of Wigner matrices I}
\author{Debapratim Banerjee}
\date{Source: arXiv:2503.23940v2 (CC BY 4.0)}
\begin{document}
\maketitle

\noindent\emph{Source and licence.} Operator limit of Wigner matrices I. Version arXiv:2503.23940v2 (CC BY 4.0), distributed under the Creative Commons Attribution 4.0 International licence, as recorded in the arXiv licence metadata for this exact version and shown on the abstract page https://arxiv.org/abs/2503.23940v2. Full rights notice: licenses/arxiv-2503.23940-CC-BY-4.0.txt.\par\smallskip

\noindent\textit{Editorial scope.} Counted unit: the Lemma whose source label is \texttt{None} in the article (source0.tex lines 471--473, source label \texttt{None}), delivered together with the article's own complete original proof of it (source0.tex lines 474--516, one \texttt{proof} environment of 43 lines). No mathematics is written, repaired or re-derived: statement and proof are the article's own text line for line. Required context retained uncounted: eq:bound (lines 440--449). Every definition, display and label of those lines is retained unchanged. Omitted from this package and not redistributed: 1--439, 450--470, 517--928. Nothing inside a retained excerpt is omitted or altered: every retained body is delivered exactly as the article's lines are written, so no omission receipt of any kind accompanies this package. Citations: the retained excerpts carry no citation command at all, so no bibliography entry of the article is transcribed and no reference is redistributed. Reference adaptation: none was needed and none was made; every reference inside the delivered unit resolves inside it, so no unresolved reference or citation is produced. Printed slips kept literal: no printed word, symbol, hypothesis or number of the counted statement or of its proof was altered. Layout and class: the article's own class is \texttt{amsart} with packages including enumerate, geometry, amsmath, amssymb, amsthm, dsfont, graphicx, microtype, xcolor, hyperref; the wrapper is the lane's pinned \texttt{amsart} editorial wrapper at 10pt on A4, which restates the theorem-like environments the retained text needs and adds the article's own macro definitions that the retained text needs (\texttt{\textbackslash tilde}), with extra wrapper packages (bm, bbm, dsfont, xspace, cancel, mathtools). The delivered numbering is the unit's own. Assets: no retained span contains a figure environment, a graphics-inclusion command, a PDF-page inclusion, a table or a TikZ picture; no image file is copied into this package, the package declares no asset dependency and no image file is redistributed with it. Licence: version arXiv:2503.23940v2 of this article is distributed under the Creative Commons Attribution 4.0 International licence, as recorded in the arXiv licence metadata for this exact version and shown on the abstract page https://arxiv.org/abs/2503.23940v2. A full rights notice is delivered at licenses/arxiv-2503.23940-CC-BY-4.0.txt. Characters: every retained excerpt is pure ASCII, so the delivered engine, XeLaTeX with its default Latin Modern text font, reports no missing character.
    \begin{equation}\label{eq:bound}
        \begin{split}
         &\left| \sum_{j=1}^{l} \tilde{f}((a_{j-1},a_{j}]) \tilde{g}((a_{j-1},a_{j}]) (a_{j}-a_{j-1})\right|\\
         &= \left| \sum_{j=1}^{l} \tilde{f}((a_{j-1},a_{j}]) \tilde{g}((a_{j-1},a_{j}]) \sqrt{(a_{j}-a_{j-1})} \sqrt{(a_{j}-a_{j-1})}\right|\\
         & \le \left| \sum_{j=1}^{l} \tilde{f}^2((a_{j-1},a_{j}])(a_{j}-a_{j-1}) \right|^{\frac{1}{2}}\times \left| \sum_{j=1}^{l} \tilde{g}^2((a_{j-1},a_{j}])(a_{j}-a_{j-1}) \right|^{\frac{1}{2}} \\
         &= \left| \sum_{j=1}^{l} \frac{\left(\int_{a_{j-1}}^{a_{j}} f(x) dx\right)^2}{(a_{j}-a_{j-1})}\right|^{\frac{1}{2}} \times \left| \sum_{j=1}^{l} \frac{\left(\int_{a_{j-1}}^{a_{j}} g(x) dx\right)^2}{(a_{j}-a_{j-1})}\right|^{\frac{1}{2}}\\
         &\le \left| \sum_{j=1}^{l} \int_{a_{j-1}}^{a_{j}} f^{2}(x)dx \right|^{\frac{1}{2}} \times \left| \sum_{j=1}^{l} \int_{a_{j-1}}^{a_{j}} g^{2}(x)dx \right|^{\frac{1}{2}}\\
         & = ||f||_{2}||g||_{2}.
        \end{split}
    \end{equation}

\begin{lemma}
 $\sup_{f \in L^{2}[0,1] ~|~ ||f||_{2}=1}\left|\langle f, \frac{dB(x)}{\sqrt{dx}} \rangle\right|$ is well defined and is $0$ with probability $1$. Further, $\int_{0}^{1} \left( \frac{dB(x)}{\sqrt{dx}} \right)^{2} dx=1$ with probability $1$.
\end{lemma}

\begin{proof}
We at first prove the second part. We take a partition $\mathcal{P}=(0=a_{0}<a_{1}<\ldots< a_{l}=1)$ and write the Riemann type sum for the integral $\int_{0}^{1} \left( \frac{dB(x)}{\sqrt{dx}} \right)^{2} dx$. It is given by 
\begin{equation}\label{eq:dbx2}
\begin{split}
&\sum_{j=1}^{l} \left(\frac{B(a_{j})-B(a_{j-1})}{(a_{j}-a_{j-1})^{\frac{1}{2}}}\right)^{2}(a_{j}-a_{j-1})\\
&= \sum_{j=1}^{l} \left(B(a_{j}) - B(a_{j-1}) \right)^2.
\end{split}
\end{equation}
It is a well-known fact in Ito calculus that the second display in \eqref{eq:dbx2} goes to $1$ in probability.

\noindent 
Now, we move to the first part of the proof. Firstly, observe that for any partition $\mathcal{P}$,
\begin{equation}
\begin{split}
&\sum_{j=1}^{l} \left(f((a_{j-1},a_{j}])+g((a_{j-1},a_{j}])\right) \left( \frac{dB}{\sqrt{dx}} \right)\left((a_{j-1},a_{j}]\right)(a_{j}-a_{j-1})\\
&= \sum_{j=1}^{l}\left( f((a_{j-1},a_{j}]) \right)\left( \frac{dB}{\sqrt{dx}} \right)\left((a_{j-1},a_{j}]\right)(a_{j}-a_{j-1})+\\
&~~~~~~~ \sum_{j=1}^{l} \left( g((a_{j-1},a_{j}]) \right)\left( \frac{dB}{\sqrt{dx}} \right)\left((a_{j-1},a_{j}]\right)(a_{j}-a_{j-1}).
\end{split}
\end{equation}
Hence, provided the limits exist, $\langle f+g , \frac{dB}{\sqrt{dx}} \rangle = \langle f,  \frac{dB}{\sqrt{dx}}  \rangle + \langle g, \frac{dB}{\sqrt{dx}}  \rangle$. Secondly, we know by Cauchy-Schwarz inequality and \eqref{eq:bound}, 
\begin{equation}
\begin{split}
&\left|\sum_{j=1}^{l}\left( f((a_{j-1},a_{j}]) \right)\left( \frac{dB}{\sqrt{dx}} \right)\left((a_{j-1},a_{j}]\right)(a_{j}-a_{j-1})\right|^2 \\
& \le \left(\int_{0}^{1} f^{2}(x)dx \right)\times \left( \sum_{j=1}^{l} \left(B(a_{j})-B(a_{j-1})\right)^2 \right).
\end{split}
\end{equation}
Hence, when $||f_{n}|| \to 0 $, we have $\langle f_{n}, \frac{dB}{\sqrt{dx}} \rangle  \to 0$.
Now, we consider any Lebesgue measurable function $f(x)$ with $||f(x)||_{2}=1$. From the definition of measurability we know that for every $\varepsilon >0$, there is a function $g_{\varepsilon}(x)= \sum_{i=1}^{K} c_{i} \mathbb{I}_{(a_{i},b_{i}]}$ 
with $c_{i} \in \mathbb{Q}$ and $a_{i}< b_{i} \in \mathbb{Q}$ such that $|| f(x)- g_{\varepsilon}(x) ||^2 < \varepsilon$. Now, we prove $\sup_{c_{i},a_{i},b_{i} \in \mathbb{Q}}\left|\langle g_{\varepsilon}(x) , \frac{dB}{\sqrt{dx}}\rangle\right|=0$ almost everywhere.
To prove this, it is enough to consider $f$ which is equal to $c \in \mathbb{Q}$ on an interval $(a,b]$ such that $a,b \in \mathbb{Q}$ and $0$ otherwise.
Now, for any partition $\mathcal{P}$ of $(a,b]$, consider
\begin{equation}
\begin{split}
&\sum_{j=1}^{l} c \frac{B(a_{j})-B(a_{j-1})}{\left(a_{j}-a_{j-1}\right)^{\frac{1}{2}}}(a_{j}-a_{j-1})\\
&= c \times G(\mathcal{P})
\end{split}
\end{equation}
where $G(\mathcal{P}) \sim N(0, \sum_{j=1}^{l}(a_{j}-a_{j-1})^2)$. Observe that 
\[
\sum_{j=1}^{l} (a_{j}-a_{j-1})^2\le ||\mathcal{P}||(b-a) \to 0
\] 
as $||\mathcal{P}|| \to 0$. So, $\langle f, \frac{dB}{\sqrt{dx}} \rangle=0$ for any $f$ which takes value $c$ on $(a,b]$ and $0$ otherwise. Here $c,a,b \in \mathbb{Q}$. However, there are only countably many choices for $c,a,b$. Hence, we can say $\sup_{c,a,b \in \mathbb{Q}} \left|\langle f, \frac{dB}{\sqrt{dx}} \rangle\right|=0$ with probability $1$. This will imply $\sup_{c_{i},a_{i},b_{i} \in \mathbb{Q}}\left|\langle g_{\varepsilon}(x), \frac{dB(x)}{\sqrt{dx}} \rangle\right| =0$ almost everywhere. Hence, $\sup_{f \in L^{2}~|~ ||f||_{1}=1}\left|\langle f(x), \frac{dB(x)}{\sqrt{dx}} \rangle\right| =0$ almost everywhere for all $f$ measurable. 
\end{proof} 

\end{document}
